\documentclass[12pt]{article}
\usepackage[margin=1in]{geometry}
\usepackage[utf8]{inputenc}
\usepackage[english]{babel}
\usepackage{setspace}
\usepackage{url}
\usepackage{footmisc}
\usepackage{fancyhdr}

\doublespacing

\pagestyle{fancy}
\fancyhf{}
\fancyhead[R]{Vaught \thepage}

\begin{document}

\begin{center}
\Large \textbf{Classical vs. Modern Hippocratic Oath}
\end{center}

Since antiquity, physicians have recited the \textit{Hippocratic Oath} as a pledge to uphold ethical standards.\footnote{Hulkower, ``History of the Hippocratic Oath.''} 
In 1964, Dr.\ Louis Lasagna introduced a modern version that highlights patient autonomy, public health obligations, 
and the importance of scientific progress.\footnote{Lasagna, \textit{Hippocratic Oath}.} 
While many argue that this updated formulation is more suited to contemporary medicine, 
the \textit{Classical Hippocratic Oath}, attributed to Hippocrates and translated by Ludwig Edelstein,\footnote{Edelstein, \textit{Hippocratic Oath}.} 
remains strikingly relevant. Its principles of nonmaleficence, humility, and respect for professional boundaries
readily adapt to modern scientific norms. By emphasizing the careful use of expertise and a commitment to avoiding harm, 
the classical oath effectively underpins the ethical and scientific foundations of current medical practice.

\bigskip

Critics often dismiss the classical oath’s invocation of Greek gods - Apollo, Asclepius, Hygieia, and Panacea - as outdated religiosity. 
Yet this element can be read symbolically, underscoring the seriousness of a physician’s calling and the need for moral humility. 
Equally disputed is the line forbidding the administration of a ``deadly drug,''\footnote{Edelstein, \textit{Hippocratic Oath}.} 
interpreted today as a broad ban on euthanasia. Historically, however, this likely served to protect healers from abuses of power 
and from being coerced into harmful acts.\footnote{Baker Institute, ``Deciphering the Hippocratic Oath.''} 
Rather than precluding compassion for the terminally ill, it reinforces a guiding principle of ``do no harm,'' 
a precept as crucial now as in antiquity. Another point of contention is the oath’s apparent prohibition on performing surgery. 
In fact, the phrase advising doctors to ``cede this work to specialists''\footnote{Edelstein, \textit{Hippocratic Oath}.} 
simply reflects the ancient recognition that not all physicians possessed the same skill set. 
Far from hampering innovation, it upholds patient welfare by discouraging practitioners from operating outside their expertise. 
In today’s multidisciplinary healthcare environment, this aligns with the modern reliance on teams of specialists, 
each of whom brings distinct skills to patient care.\footnote{Mushtque Shah et al., ``History of Surgery in Ancient, Middle and Modern Era.''}

\bigskip

Lasagna’s modern oath is often lauded for explicitly addressing patient autonomy, technology’s ethical implications, 
and social determinants of health.\footnote{Lasagna, \textit{Hippocratic Oath}.} 
These details resonate in an era of rapid scientific progress, where genome editing, digital health technologies, 
and novel pharmaceuticals raise complex moral issues. While the classical oath does not name these innovations, 
it implicitly speaks to them through broad principles like ``benefit the sick according to [one’s] ability and judgment.'' 
Physicians must remain current with scientific advances and ensure that interventions serve patient welfare, 
a goal shared by both the classical and modern versions. Furthermore, the original oath’s emphasis on justice
- the commitment to avoid exploiting patients or providing substandard care - mirrors contemporary ideals of equity and professional integrity. 
Though modern practitioners often rely on additional guidelines (institutional review boards, informed consent protocols, etc.), 
the classical oath provides a foundation for ethical discernment, 
encouraging doctors to adapt responsibly to new procedures and technologies without losing sight of their core mission.

\bigskip

Robert K.\ Merton articulated four norms underpinning rigorous science: universalism, communalism, disinterestedness, and organized skepticism.\footnote{Merton, ``A Note on Science and Democracy.''} These concepts resonate powerfully with the \textit{Classical Hippocratic Oath}. In terms of \textit{universalism}, the oath’s condemnation of injustice implies the physician’s duty to treat all patients without discrimination, reflecting the broader scientific ideal that truth-seeking and ethical obligations know no social or cultural boundaries. \textit{Communalism} is similarly embedded in the oath’s insistence on professional cooperation and mentorship, echoing the practice of sharing research findings and expertise for the common good of both the medical community and society at large. \textit{Disinterestedness} is evident in the traditional mandate that physicians always place patient welfare above financial incentives or personal ambition, thus aligning with science’s commitment to objectivity and impartial inquiry. Finally, \textit{organized skepticism} underscores the need for humility before the complexities of human health and the importance of respecting specialized knowledge. By advising doctors to cede surgical procedures to more experienced colleagues, the oath exhibits a systemic caution similar to the peer-review process in science, which demands thorough scrutiny before new discoveries are accepted. Although the modern oath references contemporary ethical challenges more explicitly, the classical version encodes these scientific norms just as robustly, ensuring that physicians exercise rigorous, patient-centered judgment in an ever-evolving clinical landscape.


\bigskip

Both the classical and modern Hippocratic oaths serve the overarching goal of guiding physicians toward ethical and compassionate practice. 
The modern text explicitly tackles autonomy, public health, and technological advances. 
Yet, the classical oath, far from being a relic, enshrines timeless principles that remain strikingly relevant 
to today’s complex and fast-paced medical environment. 
Through its grounding in harm avoidance, dedication to expertise, and respect for professional boundaries, 
the classical version not only complements scientific rigor but also nurtures the empathy and humility essential to healing. 
Even in an age of genomic breakthroughs and electronic health records, 
the enduring wisdom of the \textit{Classical Hippocratic Oath} continues to illuminate the physician’s moral path.

\newpage
\begin{thebibliography}{9}

\bibitem{baker}
Baker Institute. ``Deciphering the Hippocratic Oath in the 21st Century.'' Accessed January 20, 2025. 
\url{https://www.bakerinstitute.org/research/why-hippocratic-oath-still-relevant}

\bibitem{edelstein}
Edelstein, Ludwig, trans. \textit{Hippocratic Oath}. Baltimore: Johns Hopkins University Press, 1943. 
Accessed January 20, 2025. 
\url{https://www.pbs.org/wgbh/nova/article/hippocratic-oath-today/}

\bibitem{hulkower}
Hulkower, Raphael. ``The History of the Hippocratic Oath: Outdated, Inauthentic, and Yet Still Relevant.'' 
\textit{Albert Einstein College of Medicine} (Bronx, New York 10461).  
\url{https://einsteinmed.edu/uploadedFiles/EJBM/page41_page44.pdf} 

\bibitem{lasagna}
Lasagna, Louis. ``Hippocratic Oath.'' Boston: Tufts University School of Medicine, 1964. 
Accessed January 20, 2025. 
\url{https://www.pbs.org/wgbh/nova/article/hippocratic-oath-today/}

\bibitem{merton}
Merton, Robert K. ``A Note on Science and Democracy.'' \textit{Journal of Legal and Political Sociology} 1 (1942): 115--126.

\bibitem{mushtque}
Mushtque Shah, Wasim Shah, et al. ``History of Surgery in Ancient, Middle and Modern Era -- A Review Article.'' 
\textit{International Journal of Creative Research Thoughts (IJCRT)} 9, no. 11 (2021): 254--258. 
\url{https://ijcrt.org/papers/IJCRT2111305.pdf}

\end{thebibliography}

\end{document}
